\documentclass[9pt,a4paper]{article}

\usepackage[T1]{fontenc}
\usepackage[utf8]{inputenc}
\usepackage[cyr]{aeguill}
\usepackage[francais]{babel}

% %\renewcommand{\baselinestretch}{1.5}  % line spacing : 1,5
% \setlength{\oddsidemargin}{2,5cm}	% Marge gauche sur pages impaires
% \setlength{\evensidemargin}{2,5cm}	% Marge gauche sur pages paires
% \setlength{\textwidth}{16cm}	% Largeur de la zone de texte (17cm)
% \setlength{\topmargin}{2,5cm}	% Pas de marge en haut

% Les packages
%\usepackage[french]{babel}
\usepackage{fancybox}
\usepackage{graphics}
\usepackage{color}
\usepackage{latexsym}
\usepackage{amsmath,amsthm,amssymb}
\usepackage[plain]{fullpage} 
\usepackage{verbatim}
\usepackage{times,epsfig}
\usepackage{listings}
\usepackage{multirow}
\usepackage{amsmath,amsthm,amssymb,amscd,txfonts,bbm} % RAJOUT ICI !!!
\usepackage{graphics,epsfig} % RAJOUT ICI !!!
%\usepackage{algorithmic}
\usepackage[final]{pdfpages} 
\usepackage{fancyvrb}
\usepackage{enumerate}
\usepackage{listingsutf8}
\usepackage{tikz}
\usetikzlibrary{arrows,shapes,snakes,automata,backgrounds,petri,positioning}
\usepackage{tikz-qtree}
\usepackage{comment}


\newenvironment{qv}
{\quote\Verbatim}
{\endquote\endVerbatim}

\usepackage{lastpage}

\newcommand{\terme}[1]{\texttt{#1}}
\newcommand{\termSet}[1]{\{\texttt{#1}\}}
\newcommand{\guill}[1]{<<\,#1\,>>}
% Macro pour l'inclusion de figure  
\newcommand{\fig}[2]{%
  \begin{figure*}[hbt]
   \begin{center}
      \includegraphics[width=6cm]{#1}
  \caption{\label{#1}  #2}
  \end{center}
 \end{figure*}
}
\newcommand{\figTex}[2]{%
  \begin{figure*}[hbt]
 \centerline{\input{#1.pstex_t}}
  \caption{\label{#1}  #2}
 \end{figure*}
}

% Macro commentaires
\definecolor{turquoise}{rgb}{0.20 ,0.60, 0.60}
\definecolor{vert}{rgb}{0.40,  0.60 ,0.00}
\definecolor{violet}{rgb}{0.80 , 0.00 , 0.80}
\definecolor{bleu}{rgb}{0.20,  0.20,  0.80}
\definecolor{rose}{rgb}{1.00,  0.20,  0.40}
\def\remVG#1{\textcolor{vert}{\texttt{Virginie: #1}}}
\def\comVG#1{\begin{footnotesize}\textcolor{vert}{\texttt{Precision de Virginie: #1}}\end{footnotesize}}

\lstset{basicstyle=\footnotesize,showstringspaces=false,language=XML,breaklines=false,columns=fullflexible,frame=shadowbox, captionpos=b}

\newenvironment{solution}[0]{
~\\ \color{red} \textbf{Solution :}\color{black}}{%
}
%\excludecomment{solution}


\parindent=0pt           
                             
\usepackage{url}

\newtheorem{Exercice}{Definition}

\lstset{language=XML,basicstyle=\footnotesize,
numberstyle=\footnotesize, 
breaklines=true
showspaces=false,         
showstringspaces=false,
showtabs=false,
frame=single,
tabsize=4}   


\begin{document}

\title{Examen NFE204 - Session 1}
\date{7 f\'evrier 2012}

\maketitle

\begin{center}
{\bf \Large Consignes}\\
\begin{tabular}{| p{15cm} |}
\hline
\\
Dur\'ee de l'examen : 2h30\\
Documents non autoris\'es\\
Calculatrice autoris\'ee\\
% Il aurait fallu les pr�venir avant
%\textbf{Les transparents imprim\'es du cours sont autoris\'es. (\textsc{VT :
%mieux pour XML en raison de la syntaxe compliqu\'ee}, \`a discuter)}\\
%Tous les documents au format papier sont autoris\'es. \\ 
Les exercices sont ind\'ependants les uns des autres. Merci de soigner votre \'ecriture.\\
Ce sujet contient \pageref{LastPage} pages.\\


\textit{Important}\,: nous n'évaluons pas les réponses par la syntaxe. ne vous
inquiétez pas d'utiliser un point-virgule à la place dun point d'exclamation ou
autres détails mineurs
\\
\hline
\end{tabular}
\end{center}

\pagebreak

\section*{Première partie : XML}
. 
\textbf{\Large Exercice 1 - XPath, XQuery, XSLT}

Voici le contenu du fichier \textit{my\_playList.xml} :
\vspace{-0.05cm} 
\lstinputlisting[inputencoding=utf8/latin1,basicstyle=\small,
language=XML,breaklines=true, caption={Fichier
my\_playList.xml}]{my_playList.xml}

\subsubsection*{XPath}
1- Quel est le r\'esultat de l'\'evaluation de chacune de ces requêtes XPath
 sur le fichier \textit{my\_playList.xml} ?
\begin{enumerate}
\vspace{-0.1cm}
\item \verb7/favoris/album/titre7
\vspace{-0.1cm}
\item \verb7/favoris/album/piste[@num="8"]7
\end{enumerate}

2- Donner une requête XPath permettant de r\'ecup\'erer les noms
les titres des pistes num\'ero 1 sur les albums (information se trouvant dans
l'attribut \texttt{num}).

3- Combien d'\'el\'ements sont contenus dans le r\'esultat de
la requ\^ete \verb7//titre/parent::*7 \'evalu\'ee sur le fichier
\textit{my\_playList.xml}? Justifiez bri\`evement votre r\'eponse.

\subsubsection*{XQuery}
4- Donner une requête XQuery utilisant la syntaxe FLOWR
permettant de r\'ecup\'erer les titres des albums, ordonn\'ees par ordre alphab\'etique,
pour lesquels toutes les pistes ont un num\'ero et
un titre. {\it NB : les conditions de s\'election de la requête ne doivent être
plac\'ees ni dans la clause For, ni dans la clause Return.}

\subsubsection*{XSLT}
5- Quel est le code XML issu de la transformation du fichier
\textit{my\_playList.xml} avec la feuille de style XSLT suivante :

\begin{center}
\begin{minipage}{\textwidth}
\lstinputlisting[inputencoding=utf8/latin1,basicstyle=\small,
language=XSLT,breaklines=true, caption={Feuille de style XSLT}]{feuille.xsl}
\end{minipage} 
\end{center}

Une indentation facilitant la lecture de votre r\'eponse sera appr\'eci\'ee.

\medskip
{\it
Bar\^eme indicatif (pour l'exercice 1) : env. 7 points\\
Temps indicatif : env. 25 mn
} 


\subsection*{Exercice 2 - les DTD}

1- Voici deux arbres de donn\'ees semi-structur\'ees (repr\'esentant les
contenus de deux fichiers XML) :\\

\begin{tabular}{| c | c |}
\hline
&\\
\begin{minipage}{0.5 \textwidth}
\begin{center} 
\begin{tikzpicture}[auto]
\footnotesize
  \tikzstyle{node}=[rounded corners=1mm,thick,draw=black!75,text
  centered,text width=1.1cm,minimum height=0.3cm] 
  \tikzstyle{data}=[rounded corners=0mm,thick,draw=black!75,text
  centered,text width=1.6cm] 
  \tikzstyle{every label}=[black]

\node[node](r) at (3,4) {favoris};
\node[node](a1) at (1,3) {album};
\draw[-] (r) -- (a1);
\node[node](a2) at (4,3) {album};
\draw[-] (r) -- (a2);
\node[node](a3) at (6,3) {album};
\draw[-] (r) -- (a3);

\node[node](i1) at (0,2) {interprete};
\draw[-] (a1) -- (i1);
\node[node](t1) at (2,2) {titre};
\draw[-] (a1) -- (t1);

\node[node](t2) at (4,2) {titre};
\draw[-] (a2) -- (t2);

\node[node](i3) at (6,2) {interprete};
\draw[-] (a3) -- (i3);

\node[data](id1) at (0,1) {N. Cole};
\draw[-] (i1) -- (id1);
\node[data](td1) at (2,1) {Unforgettable};
\draw[-] (t1) -- (td1);

\node[data](td2) at (4,1) {The Wall};
\draw[-] (t2) -- (td2);

\node[data](id3) at (6,1) {The Carpenters};
\draw[-] (i3) -- (id3);

\end{tikzpicture}
\end{center}
\end{minipage}
&
\begin{minipage}{0.5 \textwidth}
\begin{center}
\begin{tikzpicture}[auto]
\footnotesize
  \tikzstyle{node}=[rounded corners=1mm,thick,draw=black!75,text
  centered,text width=1.1cm,minimum height=0.3cm] 
  \tikzstyle{data}=[rounded corners=0mm,thick,draw=black!75,text
  centered,text width=1.6cm] 
  \tikzstyle{every label}=[black]

\node[node](r) at (3,4) {favoris};
\node[node](a1) at (1,3) {album};
\draw[-] (r) -- (a1);
%\node[node](a2) at (3,3) {album};
%\draw[-] (r) -- (a2);
\node[node](a3) at (5,3) {album};
\draw[-] (r) -- (a3);

\node[node](i1) at (0,2) {interprete};
\draw[-] (a1) -- (i1);
\node[node](t1) at (2,2) {titre};
\draw[-] (a1) -- (t1);

\node[node](i3) at (4,2) {interprete};
\draw[-] (a3) -- (i3);
\node[node](t3) at (6,2) {titre};
\draw[-] (a3) -- (t3);

\node[data](id1) at (0,1) {M. le Forestier};
\draw[-] (i1) -- (id1);
\node[data](td1) at (2,1) {Plut\^ot guitare};
\draw[-] (t1) -- (td1);

\node[data](id3) at (4,1) {Billy Joel};
\draw[-] (i3) -- (id3);

\node[data](td3) at (6,1) {The Stranger};
\draw[-] (t3) -- (td3);

\end{tikzpicture}
\end{center}
\end{minipage}
\\
&\\
{\it Arbre 1} & {\it Arbre 2}\\
\hline
\end{tabular}


\bigskip

a- Proposez une DTD que {\it Arbre 1} v\'erifie mais pas {\it Arbre 2}.\\
b- Proposez une DTD que {\it Arbre 2} v\'erifie mais pas {\it Arbre 1}.\\
c- Proposez une DTD que {\it Arbre 1} \underline{et} {\it Arbre 2} v\'erifient.
\\

{\tt Rappel : Les cardinalit\'es dans une DTD se note de la sorte :
\begin{itemize}
 	\item[(vide)] 	: exactement 1 occurence
	\item[*] 		: 0 \`a infini occurences
	\item[+] 		: 1 \`a infini occurences
	\item[?] 		: 0 ou 1 occurence
\end{itemize}}
~\\

{\it
Bar\^eme indicatif : env. 3 points\\
Temps indicatif : env. 20 mn\\
} 


\section*{Seconde partie : Gestion de donn\'ees à grande \'echelle}



\subsection*{Exercice 3 - Recherche plein texte}

{\it
Bar\^eme indicatif : env. 5 points\\
Temps indicatif : env. 20 mn\\
} 

Nous souhaitons indexer les 6 documents de notre collection.
Le tableau ci-dessous donne le nombre d'occurences de chaque terme dans le
document correspondant (par souci de visibilit\'e, nous n'avons gard\'e que les mots pertinents):\\
\begin{tabular}{*{7}{|l}|}
\hline
Mot	&Doc1	&Doc2	&Doc3	&Doc4	&Doc5	&Doc6\\
\hline
\hline
Ecole	&5	&4	&2	&0	&1	&3\\
\hline
Ing\'enieur	&2	&5	&3	&3	&0	&3\\
\hline
Paris	&0	&0	&1	&0	&3	&2\\
\hline
CNAM	&0	&5	&1	&4	&3	&3\\
\hline
El\`eve	&3	&1	&1	&0	&0	&2\\
\hline
\end{tabular}

\begin{enumerate}
  \item Cr\'eer les listes inverses avec pour chaque mot\,:
  	\begin{itemize}
  	  \item La valeur de l'{\tt idf}\,;
  	  \item La liste des valeurs {\tt tf} pour chaque document correspondant.
   	\end{itemize} 
\begin{solution}

\begin{tabular}{*{8}{|l}|}
\hline
Mot		& Idf		&Doc1	& Doc2	&Doc3	&Doc4	&Doc5	&Doc6\\
\hline
\hline
Ecole	& 0,079 	&1/2	&4/15	&2/8	&		&1/7	&3/13\\
\hline
Ing\'enieur	&0,079	&1/5	&5/15	&3/8	&3/7	&		&3/13\\
\hline
Paris	&0,301		&		&		&1/8	&		&3/7	&2/13\\
\hline
CNAM	&0,079		&		&5/15	&1/8	&4/7	&3/7	&3/13\\
\hline
El\`eve	&0,176		&3/10	&1/15	&1/8	&		&		&2/13\\
\hline
\end{tabular}

\end{solution}

	\item Soit la requ\^ete suivante : ``Ecole $\land$ Paris $\land$ El\`eve"\\
	Donner les r\'esultats par ordre de pertinence norm\'ee (avec leur score correspondant) 
\begin{solution}

\begin{tabular}{*{8}{|l}|}
\hline
Documents	&Doc1	&Doc2	&Doc3	&Doc4	&Doc5	&Doc6\\
\hline
Norme		&0,068	&0,044	&0,057	&0,056	&0,133	&0,062\\
\hline
\end{tabular}

\begin{itemize}
  \item[Doc1] $ \frac{\frac{1}{2} \times 0,079 	+ 0 \times 0,301 			+ \frac{3}{10} \times 0,176}{0,068} = 0,881$
  \item[Doc2] $ \frac{\frac{4}{15} \times 0,079 + 0 \times 0,301 			+ \frac{1}{15} \times 0,176}{0,044} = 0,581$
  \item[Doc3] $ \frac{\frac{2}{8} \times 0,079 	+ \frac{1}{8} \times 0,301	+ \frac{1}{8} \times 0,176}{0,057} 	= 1,156$
  \item[Doc4] $ \frac{0 \times 0,079 			+ 0 \times 0,301 			+ 0 \times 0,176}{0,056}			= 0$
  \item[Doc5] $ \frac{\frac{1}{7} \times 0,079 	+ \frac{3}{7} \times 0,301 	+ 0 \times 0,176}{0,133}			= 1,054$
  \item[Doc6] $ \frac{\frac{3}{13} \times 0,079 + \frac{2}{13} \times 0,301 + \frac{2}{13} \times 0,176}{0,062} = 1,210$
\end{itemize}
~\\
R\'esultat tri\'e par ordre de pertinence d\'ecroissant : Doc6, Doc3, Doc5, Doc1, Doc2
\end{solution}
\end{enumerate}

\subsection*{Exercice 4 - Indexation distribu\'ee}



 {\it
Bar\^eme indicatif : env. 5 points\\
Temps indicatif : env. 20 mn\\
} 

La figure~\ref{chord-friends} montre un anneau de distribution de type
\textit{consistent hashing}. Les ronds oranges représentent des serveurs.
On veut gérer une collection de documents $D$, chacun doté d'un identifiant
$id$.


 \fig{chord-friends}{Situation initiale de l'anneau \textit{Consistent hashing}}
 
\begin{enumerate}
  \item Donnez une règle de distribution permettant d'affecter chaque document
  de $D$ à un serveur. Indiquez, selon cette règle, où se trouvent les documents
  $id=13$ et $id=16$. 
  \item On ajoute un serveur dont la valeur de hachage vaut 9. Où se place-t-il
  dans l'anneau, et quelles sont les opérations à effectuer à son arrivée\,?
  \end{enumerate}
  
\begin{solution}
 (1) On applique modulo(8) sur l'id, et on stocke l'objet sur le serveur qui
précède sur l'anneau (ou qui suit, cela revient au même). (2) Le serveur se
place en position 1. Tous les objets stockées entre 1 et 2, venant du serveur 7,
doivent être placés sur le nouveau serveur (dans l'hypothèse ``précède'').
\end{solution}
   On se pose la question de l'emplacement de la table de hachage (donnant
  l'équivalence entre les IP des serveurs et leur emplacement sur l'anneau).
  
\begin{enumerate}
  \item  On
  choisit initialement la solution suivante\,: chaque serveur connaît
  l'emplacement de son successeur sur l'anneau. Quel est le nombre de messages
  dans le pire des cas pour trouver l'emplacement d'un objet\,?
  \item Vous arrivez et proposez une solution de type \textit{Chord}, comme vu
  en cours. Quels sont les amis du serveur à l'emplacement 4\,?
\end{enumerate}
\begin{solution}
 (1) complexité linéaire\,: il faut parcourir tout l'anneau dans le pire des
 cas. (2) les amis sont 6 et 7. 
\end{solution}

\end{document}
